\subsection{The exact exterior transfer to the Jacquet norm}
\label{archimedean-jacquet-gram}

Retain the standard-character integer family, original measures,
period-normalized maps $B_\mu,B'_\mu$, the exact normalized-lowering
inputs, and gap hypotheses of
Section~\ref{archimedean-first-length}. Put $n=2r+k+1$ and
$m=2r+2$. The original signed exterior-Pieri decomposition
is supplied:
\[
 \bigwedge^\ell\mathbb C^n\otimes H_\mu
       =\bigoplus_{|I|=\ell}i_IH_{\mu+\mathbf1_I},\qquad
 [i_Ix,i_Iy]=\epsilon_I\langle x,y\rangle,\quad\epsilon_I=\pm1.
\]
Its summands are orthogonal for the displayed Hermitian form.
The inclusions need not be positive isometries. Write $p_I$ for
their projections, and $\omega_\ell,\omega'_\ell$ for the original
wedges of the first $\ell$ flag vectors. Their period-normalized
scalars satisfy
\[
 B_{\mu+\mathbf1_I}p_I(\omega_\ell\otimes v)=\kappa_I B_\mu v,\quad
 B_\mu(\operatorname{ev}_{\omega_\ell}i_Iy)
 =\epsilon_I\overline{\kappa_I}B_{\mu+\mathbf1_I}y,
\]
with corresponding $\kappa'_I$ on the right. The adjoint identity
uses the original period formula; equivariance alone is insufficient.
Zero shifted maps have scalar zero. All formal degrees $d_\mu$
refer to the same original group Haar measure.

Use inner products linear in the first variable. For a finite
coefficient $F(g)=\sum_a\langle\sigma_\mu(g)y_a,x_a\rangle A_a$ put
$s_\mu(F)=\sum_a\langle A_aB'_\mu y_a,B_\mu x_a\rangle$.
Let $\mathcal N(F)$ be the original radial squared
Hilbert--Schmidt norm of $JF$. The exterior Hermitian coefficient
$D_\ell(g)=\langle(\bigwedge^\ell g)\omega'_\ell,\omega_\ell\rangle$
is precisely the original leading minor, equal to
$\mathrm{Dbig}_\ell$ on the big cell.

\begin{proposition}\label{jacquet-exterior-exact-gram}
There is one $c_{\rm rad}>0$, fixed by the original measures and
compact data, such that
\[
 \langle JF,JH\rangle_{\rm rad}
 =c_{\rm rad}s_\mu(F)\overline{s_\mu(H)}/d_\mu.
\]
Mixed pairings from inequivalent parameters vanish. Moreover,
\begin{equation}\label{jacquet-exterior-positive-sum}
 \mathcal N(D_\ell F)
 =c_{\rm rad}|s_\mu(F)|^2
       \sum_{|I|=\ell}\frac{|\kappa_I\kappa'_I|^2}
                            {d_{\mu+\mathbf1_I}}.
\end{equation}
The left side equals the original integral
$\int|D_\ell|^2\|JF\|_{\rm HS}^2\,d\nu_{\rm rad}$.
\end{proposition}
\begin{proof}
Apply the existing mixed unfolding identity
(Theorem~\ref{jp-pairing-equality}) to finite coefficient sums.
Schur orthogonality contributes $d_\mu^{-1}$; the two original
subgroup period identities turn the finite double sum into
$s_\mu(F)\overline{s_\mu(H)}$. The radial integration formula
contributes one $c_{\rm rad}$, independent of $\mu$.
The original gap supplies the required coefficient decay;
the supplied absolute period integrals justify the diagonal
calculation. For inequivalent parameters the group pairing is
zero for every subgroup parameter before regularization.

Both period subgroups fix the flag wedges by unitriangularity.
Hence $J(D_\ell F)=D_\ell JF$ in the same original integral.
The signed tensor decomposition gives $D_\ell F=\sum_I F_I$,
where
$s_{\mu+\mathbf1_I}(F_I)
 =\epsilon_I\kappa'_I\overline{\kappa_I}s_\mu(F)$.
Every $F_I$ is a finite compact-type coefficient. The gap
preserves the parameter ordering, so different $I$ give
inequivalent representations. The mixed orthogonality just
proved gives \eqref{jacquet-exterior-positive-sum}.
\end{proof}

\begin{theorem}\label{jacquet-leading-minor-integer-family}
Let $0\le\ell\le r+1$ and set
$a_0=0$, $a_1=1$, $a_\ell=\ell+1$ for $\ell\ge2$.
If the supplied family has reserve
$\|\nu\|_1+\ell+a_\ell\le q_{\rm fam}$ at its starting index, then
\[
 \mathcal N(D_\ell F)
 \le\binom{n}{\ell}(TT')^\ell b^{8m^2}\mathcal N(F).
\]
For $\ell=0$ the norms agree without that reserve.
If $T,T'\le b^{Cm}$, $C\ge0$, the right side is at most
$b^{(C+9)m^2}\mathcal N(F)$.
\end{theorem}
\begin{proof}
The proved isometries
$T_\mu=\sqrt{d_\mu}B_\mu\sigma_\mu(\cdot)$ have mutually
orthogonal images: their cross pairing defines a bounded
intertwiner. The tensor identity and Cauchy--Schwarz imply
\[
 |\kappa_I|^2\frac{d_\mu}{d_{\mu+\mathbf1_I}}\|y\|^2
 \le\int\|(\bigwedge^\ell g^{-1})\omega_\ell\|^2
                   \|T_{\mu+\mathbf1_I}y(g)\|^2\,dq.
\]
To check the indefinite form, write $i_Iy=(x_e)_e$ in a
positive orthonormal exterior basis and test the tensor identity
against $T_{\mu+\mathbf1_I}y$ with input $\Gamma e\otimes x_e$.
The signed adjoint and Cauchy--Schwarz give
$|\beta_I|^2\|x_e\|^2
 \le\int|\langle\rho(g)\Gamma e,\omega_\ell\rangle_\Gamma|^2
                  \|T_{\mu+\mathbf1_I}y(g)\|^2dq$,
where $|\beta_I|^2=|\kappa_I|^2d_\mu/d_{\mu+\mathbf1_I}$.
Handle $x_e=0$ before cancellation, then sum and use
$\|y\|^2\le\sum_e\|x_e\|^2$. The scalar sum is the positive
exterior norm in the display.

Gram--Schmidt identifies it with $\prod_{a<\ell}L_a^2$.
Apply Corollary~\ref{original-flag-product-norm-bound}
to a nonzero lowest compact-type vector in the shifted family.
The vector has level zero, and the reserve covers its exterior
shift and first-length moments. Thus
\[
 |\kappa_I|^2d_\mu/d_{\mu+\mathbf1_I}\le T^\ell b^{4m^2},
 \quad
 |\kappa'_I|^2d_\mu/d_{\mu+\mathbf1_I}\le(T')^\ell b^{4m^2}.
\]
Zero shifted maps contribute zero.

The formal-degree sum is exactly
$\sum_{|I|=\ell}d_{\mu+\mathbf1_I}=\binom{n}{\ell}d_\mu$.
Indeed, for $\Delta(x)=\prod_{i<j}(x_i-x_j)$ the polynomial
$\sum_{|I|=\ell}\Delta(x+\mathbf1_I)$ is alternating and has
degree at most $\deg\Delta$. Divisibility by $\Delta$ and its
top homogeneous term prove the identity. The gap at least two
preserves the chamber, so $d_\mu=c_G\Delta(\mu)$ has the same
positive constant in all summands.
Insert these facts into \eqref{jacquet-exterior-positive-sum}.
When $s_\mu(F)=0$, $JF$ vanishes almost everywhere, so its
weighted norm is zero directly. Finally
$\binom{n}{\ell}\le n^\ell\le b^\ell$ and $2\ell\le m$ give
the stated power.
\end{proof}

The original integer reserve $(r+1)(r+2)$ suffices, including
$r=0$. We next retain the original real budget and both character
scales \cite{jacquet-gram-original}.

\begin{proposition}\label{jacquet-original-budget-transport}
Suppose the original $\mathrm{LocOk}_0$ conditions hold, with
$C_q\ge200$, $C_g\ge10^6$, and
\[
 \|\lambda-\lambda_0\|_1+(r+1)(r+2)
 \le Q:=\min\{h_{\rm NC}(\lambda_0)/(C_q\sqrt m),
                        g_{\rm cloud}(\lambda_0)\}/C_q .
\]
Then $\nu=\lambda-\lambda_0$ is integral, and the original
numerical family predicate holds on both sides with
$q_{\rm fam}=\lfloor Q\rfloor$ and with the full reserve
$\|\nu\|_1+(r+1)(r+2)\le q_{\rm fam}$.
\end{proposition}
\begin{proof}
Both parameters minus $(n-1)/2$ are integral coordinatewise.
Thus the expression on the left of the hypothesis is an integer,
and flooring loses nothing; $Q\ge q_{\rm fam}\ge2$.
For a compact type, the highest weight is $u-\rho_c$.
Its nonzero highest vector, evaluated on the $i$th diagonal
circle at angle $2\pi$, gives $u_i-(k-1)/2\in\mathbb Z$.
Hence $\lambda_{0,a}-u_i\in\mathbb Z+1/2$, as required.
The same holds for $u'$.

Interlacing makes the cloud exactly the first $k-1$ indices
for both compact types. Set $g=g_{\rm cloud}(\lambda_0)$
and $h=h_{\rm NC}(\lambda_0)$.
Retain the source's totalized finite-infimum convention.
Its empty cloud-pair set would give $Q=0$, so here both gap
sets are nonempty. Therefore
\[
 g\ge C_qQ\ge100(q_{\rm fam}+1),\qquad
 h\ge C_q^2\sqrt m\,Q\ge100(q_{\rm fam}+1).
\]
The original signature is exactly $\epsilon_{48}$; the
shape, boundedness and interlacing clauses are unchanged.
For completeness put $q=r+1$, so $n\ge2q$.
The original numerical quantities obey
$A_G\le4qn$, $d_N\le2qn$, $A_S\le2qn^2$ and $M\le2q^2$.
Consequently
\[
 2M(A_S+A_G)+A_G+3\le20q^3n^2,\qquad
 9qn^2+3A_G+3\le16qn^2.
\]
The original noncompact gap gives
\[
 \frac{h}{32\sqrt{2q}}
 \ge\frac{C_g}{2}q^4(3qn^2+n^3).
\]
This supplies both inequalities in the original
$\mathrm{GapOK}(h/4)$ gate, and $h/4\ge120n^3$.
The original shape parameter is at least $C_gn^7$.
These are all remaining clauses of the numerical family predicate.
This verifies that predicate, without asserting existence of
its representations or tensor decompositions.
\end{proof}

\begin{proposition}\label{jacquet-original-scale-transport}
For $c_\kappa=\kappa c_{\rm std}$, $\kappa>0$, and the original
period, quotient and radial measures, the estimate is
\[
 \mathcal N_\kappa(D_\ell F)
 \le\binom n\ell\,\kappa^{4S_\ell}(TT')^\ell
                   b^{8m^2}\mathcal N_\kappa(F),
 \qquad S_\ell=r\ell-\ell(\ell-1)/2.
\]
Under the original hypotheses $\kappa\in\{1,a\}$ this is at most
$b^{(2C_h+9)m^2}\mathcal N_\kappa(F)$.
Thus $C_b\ge C_h+5$ supplies the original leading-minor
norm bound, with the representation, exterior-Pieri, period-normalized
Bessel and exact normalized-lowering suppliers still explicit.
\end{proposition}
\begin{proof}
Let $t_\kappa$ be the original diagonal torus element with
first isotropic entries $\kappa^{i-r}$, middle block the identity
and reciprocal entries on the opposite flag.
It fixes both $z,z'$, centralizes their compact stabilizers,
normalizes $N$, and satisfies
$\psi_{\rm std}(n)=\psi_\kappa(t_\kappa n t_\kappa^{-1})$.
Moreover
$(\bigwedge^\ell t_\kappa^{-1})\omega_\ell
 =\kappa^{S_\ell}\omega_\ell$ on both flags.

Set $B_1=B_\kappa\sigma(t_\kappa)$, push the original $N$
measure by $\operatorname{Ad}(t_\kappa^{-1})$, and push the
quotient measure by left multiplication by $t_\kappa^{-1}$.
Keep the original group and compact Haar measures.
Substitution in the period formula gives its exact standard
normalization with these measures. The original $T,T'$ are
unchanged. Intertwining the exterior projection gives
\[
 B_{\kappa,I}p_I(\omega_\ell\otimes v)
 =B_{1,I}p_I\bigl((\bigwedge^\ell t_\kappa^{-1})\omega_\ell
                          \otimes\sigma(t_\kappa^{-1})v\bigr)
 =\kappa^{S_\ell}\kappa_I^{(1)}B_\kappa v .
\]
Thus $\kappa_I^{(\kappa)}=\kappa^{S_\ell}\kappa_I^{(1)}$,
and similarly on the right. Apply the standard moment theorem
with its exact pushed measures, and the positive Gram identity
with the original character and original radial measure.
This proves the first assertion without changing the big-cell
coordinates. The moment theorem requires $T,T'>0$, and does
not require either to be at least one.

Set $X=T(\kappa/a)^{2r}$ and $X'=T'(\kappa/a)^{2r}$.
The exact factor is
\[
 \kappa^{4S_\ell}(TT')^\ell
 =(XX')^\ell a^{4r\ell}\kappa^{-2\ell(\ell-1)}.
\]
For $\kappa=1$ the exponent of $a$ is $4r\ell$;
for $\kappa=a$ it is $2\ell(2r-\ell+1)\in[0,4r\ell]$.
The original $b>1$, $a,a^{-1}\le b^{C_h}$ imply $C_h\ge0$.
Using $X,X'\le b^{C_hm}$ and $a\le b^{C_h}$ yields
\[
 \kappa^{4S_\ell}(TT')^\ell
 \le b^{2C_hm\ell+4C_hr\ell}\le b^{2C_hm^2}.
\]
Finally $\binom n\ell\le b^\ell$ and $\ell\le m/2$ give the
claimed exponent, including $r=0$ and $\ell=0$.
\end{proof}

The representation, exterior-Pieri and period-normalized Bessel
existence, the exact normalized-lowering inputs, the weighted
orbital estimate and the global return sum remain separate.
The exact local formula does not identify the automorphic
global period Gram operator or close W08.
